\section{Results}
\label{sec:results}

\subsection{Decoding and encoding on one core}

Tables~\ref{tab:decode} and~\ref{tab:encode} sum the per-image times over the 40 images. With
AVX2, decoding is about 11 times faster than NBIS and encoding 4 to 6 times; SSE2, which every
x86-64 processor has, keeps 9 and 4 times. Even the scalar backend is more than five times
faster than NBIS at decoding: most of the gain comes from the algorithms (cached schedules,
branch-free entropy decoding, fewer passes over memory), not from the vector width.

% Generated by tools/gen_tables.py; do not edit by hand.
\begin{table}[ht]
  \centering\small
  \begin{tabular}{@{}lrrrrrr@{}}
    \toprule
    & \multicolumn{3}{c}{0.75 bits per pixel} & \multicolumn{3}{c}{2.25 bits per pixel}\\
    \cmidrule(lr){2-4}\cmidrule(l){5-7}
    implementation & total (ms) & Mpx/s & speedup & total (ms) & Mpx/s & speedup\\
    \midrule
    NBIS 5.0.0 (C) & 1\,404 & 28 & 1.0$\times$ & 1\,622 & 24 & 1.0$\times$\\
    scalar, exact & 253 & 153 & 5.5$\times$ & 284 & 137 & 5.7$\times$\\
    scalar, fast & 226 & 172 & 6.2$\times$ & 258 & 150 & 6.3$\times$\\
    SSE2, exact & 155 & 250 & 9.0$\times$ & 179 & 216 & 9.0$\times$\\
    SSE2, fast & 154 & 252 & 9.1$\times$ & 179 & 217 & 9.1$\times$\\
    AVX2, exact & 119 & 325 & 11.8$\times$ & 150 & 259 & 10.8$\times$\\
    AVX2, fast & 120 & 324 & 11.7$\times$ & 149 & 260 & 10.9$\times$\\
    \bottomrule
  \end{tabular}
  \caption{Decoding the 40 NBIS-encoded images of the NIST set (41.0 megapixels in all), one core: sum of the per-image medians, throughput and speedup over NBIS.}
  \label{tab:decode}
\end{table}

% Generated by tools/gen_tables.py; do not edit by hand.
\begin{table}[ht]
  \centering\small
  \begin{tabular}{@{}lrrrrrr@{}}
    \toprule
    & \multicolumn{3}{c}{0.75 bits per pixel} & \multicolumn{3}{c}{2.25 bits per pixel}\\
    \cmidrule(lr){2-4}\cmidrule(l){5-7}
    implementation & total (ms) & Mpx/s & speedup & total (ms) & Mpx/s & speedup\\
    \midrule
    NBIS 5.0.0 (C) & 1\,080 & 36 & 1.0$\times$ & 1\,210 & 32 & 1.0$\times$\\
    scalar, exact & 442 & 88 & 2.4$\times$ & 495 & 78 & 2.4$\times$\\
    scalar, fast & 370 & 105 & 2.9$\times$ & 425 & 91 & 2.8$\times$\\
    SSE2, exact & 278 & 140 & 3.9$\times$ & 341 & 114 & 3.6$\times$\\
    SSE2, fast & 221 & 175 & 4.9$\times$ & 288 & 135 & 4.2$\times$\\
    AVX2, exact & 242 & 160 & 4.5$\times$ & 313 & 124 & 3.9$\times$\\
    AVX2, fast & 186 & 209 & 5.8$\times$ & 251 & 155 & 4.8$\times$\\
    \bottomrule
  \end{tabular}
  \caption{Encoding the 40 images of the NIST set at two target rates, one core.}
  \label{tab:encode}
\end{table}


Table~\ref{tab:stages} breaks the times of \texttt{a002} into stages. The inverse transform
still dominates decoding (about 3.2 of 4.8~ms); entropy decoding and dequantization take
1.0~ms together. In the encoder, the forward transform takes half the time and the tokenizer is
the next largest stage.

% Generated by tools/gen_tables.py; do not edit by hand.
\begin{table}[ht]
  \centering\small
  \begin{tabular}{@{}lrrrrrrr@{}}
    \toprule
    decoder & parse & zero & entropy & dequantize & plan & transform & output\\
    \midrule
    scalar, exact & 0.03 & 0.21 & 0.62 & 0.37 & 0.01 & 7.37 & 1.13\\
    AVX2, exact & 0.03 & 0.22 & 0.63 & 0.37 & 0.01 & 3.17 & 0.34\\
    AVX2, fast & 0.03 & 0.21 & 0.62 & 0.34 & 0.01 & 2.96 & 0.31\\
    \midrule
    encoder & normalize & transform & variance & quantize & tokenize & tables & entropy\\
    \midrule
    scalar, exact & 0.47 & 8.84 & 0.19 & 2.72 & 1.57 & 0.19 & 0.39\\
    AVX2, exact & 0.43 & 3.84 & 0.18 & 0.61 & 1.27 & 0.19 & 0.39\\
    AVX2, fast & 0.42 & 3.98 & 0.19 & 0.61 & 1.30 & 0.19 & 0.39\\
    \bottomrule
  \end{tabular}
  \caption{Median time of each stage, in milliseconds, for image \texttt{a002} at 0.75 bits per pixel. Decoder: \emph{zero} clears the index buffer, \emph{plan} fetches the cached transform schedules. Encoder: \emph{entropy} writes the bit stream; table building and tokenizing are separate.}
  \label{tab:stages}
\end{table}


\subsection{Threads and batches}

A single image decodes 1.7 times faster with four threads (Table~\ref{tab:threads}); the
serial parts (the deep splits, the largest entropy-coded block) limit the gain, and small
images gain little. Batches scale better, since images are independent: on the four performance
cores, 829 images per second, against 319 on one core (Table~\ref{tab:batch}). Eight threads
(the performance cores with both their hardware threads) reach 934; twelve, adding the
efficiency cores, are slower, because the transform saturates memory bandwidth.

% Generated by tools/gen_tables.py; do not edit by hand.
\begin{table}[ht]
  \centering\small
  \begin{tabular}{@{}lrrrr@{}}
    \toprule
    image & 1 thread & 2 threads & 4 threads & 8 threads\\
    \midrule
    \texttt{cmp00001} (589$\times$605) & 0.96 & 0.95 & 0.79 & 0.84\\
    \texttt{a002} (1452$\times$1006) & 4.63 & 3.94 & 2.66 & 2.73\\
    \texttt{sample\_19} (1600$\times$1500) & 7.69 & 6.63 & 4.58 & 4.42\\
    \bottomrule
  \end{tabular}
  \caption{Decoding one image with several threads (AVX2, exact): median milliseconds. Threads beyond four share the four performance cores.}
  \label{tab:threads}
\end{table}

% Generated by tools/gen_tables.py; do not edit by hand.
\begin{table}[ht]
  \centering\small
  \begin{tabular}{@{}lrrrrr@{}}
    \toprule
    threads & 1 & 2 & 4 & 8 & 12\\
    \midrule
    exact, images/s & 319 & 530 & 829 & 934 & 843\\
    exact, Mpx/s & 309 & 513 & 803 & 904 & 817\\
    fast, images/s & 312 & 520 & 790 & 856 & 839\\
    fast, Mpx/s & 302 & 504 & 765 & 829 & 813\\
    \bottomrule
  \end{tabular}
  \caption{Batch decoding of the 40 images at 0.75 bits per pixel, one image per thread (AVX2). Without the inverse transform (entropy decoding and dequantization only), four threads reach 3\,093 images/s.}
  \label{tab:batch}
\end{table}


\subsection{WebAssembly}

The extension runs the codec as WebAssembly. With SIMD the WebAssembly build is 1.35 times
slower than native AVX2 and still eight times faster than native NBIS; without SIMD it is about
three times slower than with it (Table~\ref{tab:wasm}). Its outputs are identical to NBIS's: WebAssembly
float arithmetic follows IEEE 754 exactly, and the exact kernels use no fused operations.

% Generated by tools/gen_tables.py; do not edit by hand.
\begin{table}[ht]
  \centering\small
  \begin{tabular}{@{}lrrrr@{}}
    \toprule
    image & NBIS (C) & native AVX2 & wasm scalar & wasm SIMD\\
    \midrule
    \texttt{cmp00001} & 12.31 & 0.96 & 4.62 & 1.54\\
    \texttt{a002} & 52.38 & 4.80 & 18.67 & 6.50\\
    \texttt{sample\_19} & 86.51 & 7.65 & 29.14 & 10.74\\
    \bottomrule
  \end{tabular}
  \caption{Decoding in WebAssembly (Node.js 24, one core), exact mode, median milliseconds; the WebAssembly outputs are identical to those of NBIS.}
  \label{tab:wasm}
\end{table}


\subsection{A server processor with AVX-512}
\label{sec:server}

The laptop has no AVX-512, so the test suite and the single-core measurements were repeated on a
server processor that has it, an Intel Xeon Gold 6330, which other users shared at the time.
Every test passes there with AVX-512 among the backends: the exact decoder reproduces NBIS on the
86 decoder files with 1 and 4 threads, and the exact encoder writes NBIS's bytes in the 80 cases.

Table~\ref{tab:server} gives the times. Doubling the vector width from AVX2 to AVX-512 gains
little, as expected from a transform limited by data movement (Section~\ref{sec:optimizations}).
In four rounds that alternated the two backends on the same core, exact decoding was 8 to 15\%
faster with AVX-512 in every round and exact encoding 2 to 8\% (one round, at 27\%, was an
outlier); fast decoding and fast encoding showed no consistent difference. The absolute times
varied by up to 20\% between rounds, because of the other jobs on the server, so only the ratios
within a round are reliable. A server core is slower than the laptop's performance core (NBIS takes
2.5~s for the set against 1.4~s), but the speedup over NBIS is similar: about 12 to 13 times at
decoding and 4.5 to 5.7 times at encoding.

% Generated by tools/gen_tables.py; do not edit by hand.
\begin{table}[ht]
  \centering\small
  \begin{tabular}{@{}lrrrrrr@{}}
    \toprule
    & \multicolumn{3}{c}{decode} & \multicolumn{3}{c}{encode}\\
    \cmidrule(lr){2-4}\cmidrule(l){5-7}
    implementation & total (ms) & vs NBIS & vs AVX2 & total (ms) & vs NBIS & vs AVX2\\
    \midrule
    NBIS 5.0.0 (C) & 2\,546 & 1.0$\times$ & 0.08 & 1\,775 & 1.0$\times$ & 0.22\\
    scalar, exact & 477 & 5.3$\times$ & 0.43 & 754 & 2.4$\times$ & 0.53\\
    scalar, fast & 439 & 5.8$\times$ & 0.47 & 652 & 2.7$\times$ & 0.61\\
    SSE2, exact & 253 & 10.1$\times$ & 0.81 & 470 & 3.8$\times$ & 0.85\\
    SSE2, fast & 260 & 9.8$\times$ & 0.79 & 419 & 4.2$\times$ & 0.95\\
    AVX2, exact & 205 & 12.4$\times$ & 1.00 & 397 & 4.5$\times$ & 1.00\\
    AVX2, fast & 193 & 13.2$\times$ & 1.06 & 312 & 5.7$\times$ & 1.28\\
    AVX-512, exact & 188 & 13.5$\times$ & 1.09 & 391 & 4.5$\times$ & 1.02\\
    AVX-512, fast & 196 & 13.0$\times$ & 1.05 & 312 & 5.7$\times$ & 1.28\\
    \bottomrule
  \end{tabular}
  \caption{The same measurements on a server processor with AVX-512 (Intel Xeon Gold 6330, one core): the 40 images of the NIST set at 0.75 bits per pixel, sum of the per-image medians. \emph{vs AVX2} is the speed relative to the exact AVX2 backend on the same core.}
  \label{tab:server}
\end{table}

